\clearpage
\fancyhead[L]{\small CERES.}
\hypertarget{native-317-lower}{}
\begin{center}
\textbf{CERES.}
\par\vspace{4pt}\hrule\vspace{4pt}
\fontsize{10}{12}\selectfont Göttingische gelehrte Anzeigen. Issue 87. Pp. 857..860. 3 June 1809.
\par\vspace{4pt}\hrule\vspace{4pt}
\end{center}

\par\indent In issue 40 of last year’s Anzeigen [p. 299 of this volume] we reported the twelfth set of elements of \emph{Ceres}, which Professor \textsc{Gauss} had based on observations from 1807. We now present the results of the calculations just completed from the observations of 1808. The planet’s low position in Capricorn made it very difficult to observe in our northern regions. At our observatory only a single meridian observation could be made, and even that one cannot be considered particularly good.


\begin{center}\fontsize{10}{12}\selectfont\setlength{\tabcolsep}{5pt}\renewcommand{\arraystretch}{1.15}\begin{tabular}{lr} 1808 mean time July 25. & $13^{\text{u}}\;14^{\text{m}}\;29^{\text{s}}$.\\ Right ascension & $322^\circ\;8'\;57''6$.\\ South declination & $27^\circ\;52'\;12''5$\\
\end{tabular}\end{center}


\SourceLeaf{318}\par\indent The only other observations that have become known so far are those made by Mr. von \textsc{Lindenau} at the Seeberg Observatory and Mr. \textsc{Santini} in Padua. With these observations to work from, Professor \textsc{Gauss} determined the opposition of 1808, the error of the XII elements, and the corrections needed to bring these elements into the closest possible agreement with all six oppositions observed so far. With correspondence now so difficult, we hope astronomers will welcome the publication of these results in our journal.

\par\indent The opposition was determined by examining the observations as carefully as possible, with the following result.

\begin{center} 1808 Aug. 5. \quad $11^{\text{u}}\ 30^{\text{m}}\ 13^{\text{s}}$ mean time in Göttingen

true longitude of Ceres .\;.\;.\;\ $313^\circ\ 13'\ 40''2$

true geocentric latitude .\;.\;\ $12\ 43\ 44{,}8$ south
\end{center}

\par\indent This time the XII elements differ by only one minute in longitude and a couple of seconds in latitude. That difference is quite insignificant, considering that the perturbations used are still incomplete and that perfect agreement among all observations made over nearly eight years must in any case be regarded as impossible. Professor \textsc{Gauss} has already laid out everything needed for the future, more complete calculation of the perturbations, though he cannot yet say when he will be able to carry out this work, which is as delicate as it is extensive. In the meantime he thought it worth investigating how closely all six observed oppositions could still be reconciled using the perturbation tables he calculated in 1802. The result shows that the differences are again so small that they give no reason to include the more complete perturbations for some time yet. It should be borne in mind, however, that when the perturbations have to be taken fully into account, calculating a single planetary position will certainly take four times as much time and work as it does now. The new elements of Ceres, which differ very little from the XII elements, are given below. The epochs refer to the meridian of Göttingen.

\Needspace{23\baselineskip}
\begin{center} XIII. \emph{Elements of Ceres.}
\end{center}


\begin{center}\fontsize{10}{12}\selectfont\setlength{\tabcolsep}{5pt}\renewcommand{\arraystretch}{1.15}\begin{tabular}{crr}  & Mean longitude. & Perihelion.\\\hline 1801 & $77^\circ\;18'\;36''5$ & $146^\circ\;26'\;0''1$\\ 1802 & $155\;28\;23.4$ & $\phantom{146}\;28\;1.4$\\ 1803 & $233\;38\;10.3$ & $\phantom{146}\;30\;2.6$\\ 1804 & $312\;0\;48.1$ & $\phantom{146}\;32\;4.2$\\ 1805 & $30\;10\;35.0$ & $\phantom{146}\;34\;5.4$\\ 1806 & $108\;20\;21.8$ & $\phantom{146}\;36\;6.6$\\ 1807 & $186\;30\;8.7$ & $\phantom{146}\;38\;7.9$\\ 1808 & $264\;52\;46.5$ & $\phantom{146}\;40\;9.5$\\ 1809 & $343\;2\;33.4$ & $\phantom{146}\;42\;10.7$\\
\end{tabular}\end{center}

\begin{center}\fontsize{9}{11}\selectfont\setlength{\tabcolsep}{5pt}\renewcommand{\arraystretch}{1.15}\begin{tabular}{lr} Daily mean tropical motion .\; & $770''9230$\\ Eccentricity 1806 .\;.\;.\;.\;.\;.\;.\;.\; & $0.0785028$\\ Annual decrease .\;.\;.\;.\;.\;.\;.\;.\; & $0.00000583$\\ Logarithm of the semimajor axis .\;.\; & $0.4420486$\\ Ascending node 1806 .\;.\;.\;.\; & $80^\circ\;53'\;41''3$\\ Annual motion .\;.\;.\;.\;.\;.\;.\;.\; & $+1{,}48$\\ Orbital inclination 1806 .\;.\;,\;.\;.\; & $10^\circ\;37'\;31''2$\\ Annual decrease .\;.\;.\;.\;.\;.\;.\;.\;.\;.\; & $0{,}44$\\
\end{tabular}\end{center}


\Needspace{13\baselineskip}\par\indent The following differences show how closely these new elements agree with all six oppositions observed so far.

\SourceLeaf{319}\begin{center} Difference:
\end{center}


\begin{center}\fontsize{10}{12}\selectfont\setlength{\tabcolsep}{5pt}\renewcommand{\arraystretch}{1.15}\begin{tabular}{crr}  & \shortstack{Heliocentric longitude\\in the orbit} & \shortstack{Heliocentric\\latitude}\\\hline 1802 & $-0''5$ & $-1''9$\\ 1803 & $+22.1$ & $-5.3$\\ 1804 & $-34.9$ & $-1.2$\\ 1806 & $+26.7$ & $-3.6$\\ 1807 & $-26.9$ & $-11.7$\\ 1808 & $+20.8$ & $-6.0$\\
\end{tabular}\end{center}


\Needspace{17\baselineskip}\par\noindent It should be noted that the perturbations in longitude and in the radius vector used in this calculation come from the perturbation tables that Professor \textsc{Gauss} published in the March issue of the Monatl. Corresp. [March 1803, p. 235 and following pages of this volume]. The perturbations in latitude, by contrast, are based on a later, more complete expansion whose results are still unpublished. The small differences in latitude may be taken as confirmation that no significant terms were overlooked in that determination. [The contents of this report are reproduced in the Monatliche Correspondenz, vol. XIX, pp. 504\,..\,512, May 1809. It also prints the ephemerides calculated by \textsc{Gauss}, for Ceres from 18 July 1809 to 16 April 1810 and for Juno from 4 September 1809 to 10 May 1810. For midnight in Göttingen at four-day intervals they give geocentric right ascension and geocentric declination to the minute, and the distance from Earth to four decimal places.]

